\documentclass[11pt,UTF8,fontset=fandol]{ctexart}
\usepackage{amsmath,amssymb,amsfonts}
\usepackage{graphicx,float,caption,subcaption}
\usepackage{booktabs,array,multirow,makecell}
\usepackage{geometry,enumitem,xcolor}
\usepackage[hidelinks]{hyperref}
\usepackage{url}
\geometry{margin=2.3cm}
\graphicspath{{/Users/yye/project/论文/assets/archer/}}
\captionsetup{font=small,labelfont=bf}
\setlength{\parskip}{0.4em}
\hypersetup{pdftitle={ArCHer 中文精读导读},pdfauthor={Yifei Zhou et al.}}

% 小白知识卡片：上下加线 + 粗体标签，稳健不挑内容
\newenvironment{tip}[1]{%
  \par\smallskip\noindent\rule{\linewidth}{0.4pt}\par\nobreak
  \noindent\textbf{【小白补丁】#1}\par\nobreak\noindent\ignorespaces
}{%
  \par\noindent\rule{\linewidth}{0.4pt}\par\smallskip
}

\title{\textbf{ArCHer 中文精读导读：用“分层强化学习”教语言模型当多轮智能体}}
\author{原文：Yifei Zhou, Andrea Zanette, Jiayi Pan, Sergey Levine（UC Berkeley，加州大学伯克利分校）\\ Aviral Kumar（Google DeepMind，谷歌深度思维）}
\date{}

\begin{document}
\maketitle
\begin{center}\small\textit{本导读为面向 AI 应用开发初学者的中文学习笔记：含要点讲解、关键句翻译、公式精讲与图表解读。出于版权合规，这是转化性学习材料，以转述和讲解为主，不逐字翻译原文。}\end{center}

\section*{论文卡片}
\begin{center}
\begin{tabular}{@{}ll@{}}
\toprule
\textbf{英文标题} & ArCHer: Training Language Model Agents via Hierarchical Multi-Turn RL \\
            & （ArCHer：用分层式多轮强化学习训练语言模型智能体） \\
\textbf{作者} & Yifei Zhou, Andrea Zanette, Jiayi Pan, Sergey Levine, Aviral Kumar \\
\textbf{机构} & UC Berkeley（加州大学伯克利分校）、Google DeepMind（谷歌深度思维） \\
\textbf{会议} & ICML 2024（国际机器学习大会） \\
\textbf{引用数} & 180（Semantic Scholar，截至 2026-06-07） \\
\textbf{代码} & \url{https://github.com/YifeiZhou02/ArCHer} \\
\textbf{方向} & 多轮智能体（multi-turn agent）$\times$ 大模型 $\times$ 强化学习（RL）后训练 \\
\textbf{一句话定位} & 用\emph{分层 actor-critic}（演员-评论家）框架解决多轮长程任务里的“功劳分配”\\
              & 难题：高层用 off-policy 的 TD 学习训一个\emph{话语级}价值函数，评估“说完这句话\\
              & 局面好坏”；低层用 on-policy 的\emph{token 级}策略梯度决定这句话每个字怎么生成。\\
              & 在 Twenty Questions 等任务上，样本效率比纯 on-policy 的 PPO 约提升 100 倍。 \\
\bottomrule
\end{tabular}
\end{center}

\section{小白知识卡：读这篇前要懂的几个概念}

这篇论文站在“大模型 + 强化学习 + 智能体”的交叉点上，下面这几个词先弄明白，后面就顺了。

\begin{itemize}[leftmargin=1.5em]
\item \textbf{智能体任务（agent task）/ 多轮（multi-turn）}：让大模型不止“回答一个问题就完事”，而是要\emph{连续做好几步决策}去完成一个目标。比如：在网页上一步步操作买东西、在 Linux 终端里先读手册再敲命令、跟用户来回聊好几轮把需求问清楚。每一步叫一“轮”（turn），每一轮模型说的那句话叫一个\emph{话语}（utterance）。

\item \textbf{强化学习（reinforcement learning, RL）}：一种“试错学习”。智能体在环境里行动（action），环境给反馈（状态 state 变化 + 奖励 reward），目标是让\emph{长期累计奖励}最大。和监督学习不同，它不需要“标准答案”，只需要一个能打分的奖励信号。

\item \textbf{单轮 RL vs.\ 多轮 RL}：现在主流的 RLHF（基于人类反馈的强化学习）大多是\emph{单轮}的——给一个 prompt、生成一个回答、用奖励模型打个分、更新一下，就结束了。它不需要模型考虑“我这句话会怎么影响后面几轮”。多轮 RL 则必须考虑这种长期影响。

\item \textbf{PPO（近端策略优化，Proximal Policy Optimization）}：当前 RLHF 最常用的算法，属于 \emph{on-policy}（同策略）方法——每次更新只能用“当前这版模型刚刚跑出来的数据”，旧数据用完就扔。优点是稳，缺点是\emph{样本效率低}（很费数据/很费与环境交互的次数）。

\item \textbf{on-policy（同策略）vs.\ off-policy（异策略）}：on-policy 只能用当前策略产生的数据；off-policy 可以把\emph{过去所有交互}存进一个回放缓冲区（replay buffer）反复利用。在多轮任务里，每次与环境交互（比如真的去点一次网页）又慢又贵，所以能复用旧数据的 off-policy 方法在样本效率上有巨大优势。

\item \textbf{功劳分配 / 信用分配（credit assignment）}：长程任务的核心难题。一连串动作最后才拿到奖励，到底是\emph{哪一步}的功劳（或过错）？比如二十问游戏里赢了，是因为第 3 个问题问得好，还是第 15 个？把奖励正确地“摊派”回每一步，就是 credit assignment。

\item \textbf{actor-critic（演员-评论家）}：一类 RL 算法。\emph{actor}（演员）= 策略，负责“做什么动作”；\emph{critic}（评论家）= 价值函数，负责“评估当前局面/动作有多好”。critic 给 actor 提供更稳的学习信号，减少噪声。

\item \textbf{价值函数 / Q 函数（value / Q function）与 TD 学习}：$V(s)$ 估计“处在状态 $s$ 时，未来能拿多少累计奖励”；$Q(s,a)$ 估计“在状态 $s$ 做动作 $a$ 后能拿多少”。\emph{TD 学习}（temporal-difference，时序差分）是训练这些函数的经典办法：用“当前奖励 + 下一步的估值”来纠正“当前的估值”，即所谓 Bellman 自举（bootstrapping），不必等到一整局结束才学。

\item \textbf{token 级 vs.\ 话语级（token-level vs.\ utterance-level）}：把模型生成的\emph{每一个字（token）}都当成一个动作，叫 token 级；把\emph{一整句话}当成一个动作，叫话语级。token 级动作空间小、好优化，但一局下来 token 数量巨多、时间跨度（horizon）极长，奖励信号要传播很久；话语级时间跨度短，但“一句话”的动作空间是天文数字（所有可能的句子），TD 学习里要“在所有动作里取最大值”变得不可行。ArCHer 的巧思就是\emph{两者各取所长}。
\end{itemize}

\begin{tip}{一句话抓住这篇在干嘛}
多轮智能体训练，token 级 RL“时间跨度太长、学得太慢”，话语级 RL“动作空间太大、没法取最大值”，纯 on-policy 的 PPO 又“太费数据”。ArCHer 把它们\emph{分两层}：高层把“一句话”当一个动作，用能复用旧数据的 off-policy TD 学习训一个评委（话语级 critic）；低层把“一句话里的每个字”当动作，用 PPO 那套 on-policy 策略梯度生成，奖励就用高层评委给的分。于是同时拿到“短跨度 + 小动作空间 + 高样本效率”。
\end{tip}

\section{摘要全译}

大语言模型（large language models, LLMs）的一大类用途是面向目标的决策任务（即“智能体”任务）：模型不只是给定 prompt 生成一个像样的续写，而是要在\emph{多轮交互}中、在较长时间跨度内做出一连串聪明决策来完成任务（例如与网页交互、使用软件工具、做客服支持）。强化学习（RL）为这类智能体任务提供了通用范式，但当前面向 LLM 的 RL 方法大多聚焦于\emph{单轮}的奖励最大化。从设计上看，今天的单轮 RL 方法根本无法训练 LLM 去“跨多轮地主动搜集并整合信息”、做“功劳分配”、或对自己过去的动作进行推理——而这些恰恰是智能体任务的关键。这就提出一个问题：我们该如何为 LLM 设计\emph{有效又高效}的多轮 RL 算法？

本文提出一个用于开发“微调 LLM 的多轮 RL 算法”的算法框架，它既保留了现有单轮 RL 方法（如 PPO）的灵活性，又能有效地容纳多轮、长时间跨度和延迟奖励。为此，我们的框架采用\emph{分层 RL}（hierarchical RL）思路，并行运行两个 RL 算法：一个\emph{高层 off-policy} 的 RL 算法，训练一个价值函数来在\emph{话语}层面聚合奖励；一个\emph{低层} RL 算法，利用这个高层价值函数（取代单轮 RL 里用的奖励模型）来逐字（token-by-token）训练每一轮/每句话内部的策略。我们把这个框架命名为\textbf{ArCHer}（Actor-Critic Framework with a Hierarchical Structure，带分层结构的演员-评论家框架），它还能衍生出其他若干种 RL 方法。实验上我们发现，ArCHer 在多轮任务上显著提升了效率与性能：相比现有 on-policy 方法，样本效率提升约 \textbf{100 倍}；同时它还能很好地受益于模型容量扩大（我们测到了 70 亿参数规模）。项目主页：\url{https://yifeizhou02.github.io/archer.io/}，代码：\url{https://github.com/YifeiZhou02/ArCHer}。

\section{逐节精读}

\subsection{引言：为什么单轮 RL 教不出好智能体}

LLM 凭借通用知识，有潜力解决各种能用自然语言表达的决策/“智能体”问题：写代码、上网操作、用工具、与人交流等等。要在这些场景里成功，模型必须\emph{跨多轮做出一连串聪明决策}，而不是每步只生成“概率最高的续写”。

但现状是：大多数让 LLM 表现出目标导向行为的方法，都依赖\emph{短视的目标}——要么逐步模仿成功示范（行为克隆），要么优化单轮偏好（如 RLHF）。作者点出单轮方法的两个硬伤（这两句是全文动机的关键）：

\begin{quote}
“用单轮方法训练出来的策略，往往做不好功劳分配（例如，无法识别出那些当下看似次优、却能带来长期收益的好动作），也无法赋予策略\emph{主动搜集信息}的行为——而这在智能体问题中很重要（例如面对一个新工具时）。”
\end{quote}

\textbf{图 1 的对比直觉}：原文图 1 用一个“用户说没动力工作”的例子做对比。\emph{单轮智能体}试图在一轮里把问题解决，于是一股脑甩出一大堆建议（结果用户觉得是陈词滥调，不满意）；\emph{多轮智能体}则先反问“最近是不是发生了什么事”，先\emph{搜集信息}再针对性回应，最后用户很满意。技术上，单轮 RL 用“奖励模型 + 策略梯度”训练；本文的多轮做法是把\emph{奖励模型换成一个用 off-policy RL 学出来的价值函数}。

\textbf{多轮 RL 的独特挑战}（作者归纳）：
\begin{enumerate}[leftmargin=1.6em]
\item 多轮 RL 需要与外部环境（人、网页服务器）\emph{在线交互}，又慢又贵。因此像 PPO 这种 on-policy 方法很快变得不实用——它\emph{不能复用过去的交互数据}。
\item off-policy 的纯离线方法能复用数据，但又有新麻烦：多轮下 token 数量暴涨（LLM 还偏爱长回答），把\emph{每个 token 当动作}的方法（如 token 级 ILQL）就得在\emph{极长的时间跨度}上传播奖励，学得极慢。作者实测：token 级 ILQL 在某任务上要\emph{10 天以上}才收敛，而过滤式行为克隆不到一天。
\item 反过来，把\emph{一整句话当一个动作}能缩短时间跨度，但会带来一个\emph{巨大且变长的动作空间}——这对基于 TD 学习、每步都要“在动作上取最大值”的 off-policy 方法是灾难。
\end{enumerate}

\textbf{核心洞见}：用\emph{分层}的办法同时绕开上面所有坑。高层用 off-policy TD 学习训一个\emph{话语级}价值函数；低层用\emph{任意 on-policy 策略梯度}优化每一轮内部的 token 生成，并把高层价值函数当作“这一轮结束时的奖励”。这样既能复用样本、收敛快，又避免了“对单个 token 做 Bellman 回溯”或“在天文数字的动作空间里取最大值”。一句话总结贡献：

\begin{quote}
“这样我们就同时拿到了\emph{话语级}与 \emph{token 级}、\emph{off-policy} 与 \emph{on-policy} 各自的好处。”
\end{quote}

\subsection{相关工作：前人卡在哪}

\begin{itemize}[leftmargin=1.5em]
\item \textbf{LLM 的单轮 RL}：PPO、A2C、DPO、过滤式监督学习等，都聚焦“模型只做一个决策、不再与环境交互”的单轮偏好优化。但很多智能体问题\emph{无法在单轮内解决}——它们需要先\emph{主动搜集信息}再决策（如先问清用户偏好再订行程、先读终端帮助手册再执行命令），单轮方法学不会这种细腻策略。
\item \textbf{不用 RL 训练语言智能体}：靠提示工程（prompt engineering），如 ReAct（先“想”再做）、Reflexion（反思过去失败）、Voyager（自建课程与动作库）。但\emph{不更新模型参数}，效果天花板被预训练能力锁死；连 GPT-4 加上下文学习，在分布外场景也可能表现很差。另一条路是用\emph{成功轨迹微调}，但人工标注/工具调用标注很贵，且任务越长，自动“撞”到成功轨迹的概率越低、代价越高。
\item \textbf{LLM 的多轮 RL}：直接套用 PPO、REINFORCE 等策略梯度，在需要与环境交互的多步场景里\emph{样本效率很低}。off-policy / 离线的价值方法能用静态数据学，但现有方法要么 (1) 把\emph{单个 token 当动作}（token 级），被超长时间跨度拖累；要么 (2) 把\emph{一句话当动作}，却得从一个\emph{冻结的预训练 LLM}里采样多个候选句子来做 Bellman 回溯里的“取最大值”，这会拖慢策略改进的速度。ArCHer 要同时解决这两个局限。
\end{itemize}

\subsection{方法第一步：把语言生成建模成“分层 MDP”}

这是全文方法的地基。作者把多轮语言生成\emph{重新形式化}为在一个\emph{分层马尔可夫决策过程}（hierarchical Markov decision process, MDP）里行动——\emph{一个高层 MDP，里面嵌套一个低层 MDP}。

\begin{itemize}[leftmargin=1.5em]
\item 两层 MDP 的\textbf{状态}都是\emph{变长的 token 序列}。
\item \textbf{高层 MDP}：一个动作 $a_t$ 是\emph{一整段 token 序列（一句话/一轮）}；状态 $s_t$ 是“LLM 与环境到目前为止的交互历史”。
\item \textbf{低层 MDP}：一个动作 $a^h_t$ 是\emph{单个 token}（第 $t$ 句话里的第 $h$ 个字）。低层 MDP 建模的是“怎么把高层那一句话一个字一个字生成出来”。低层的状态 = 拼接“到这一轮为止的高层状态 $s_c$” + “本轮内已经生成的字 $a^{1:h-1}_t$”。
\item \textbf{两层如何衔接}：在低层 MDP 里执行一连串动作（吐出一串 token），\emph{恰好等于}在高层 MDP 里执行一个动作（说完一句话）。低层何时结束？当策略选择吐出 “EOS”（结束符）时。
\end{itemize}

\textbf{两层各自的目标（关键）}：高层策略最大化\emph{任务奖励}；低层策略则去找一段最多 $L$ 个 token 的序列 $a^{1:L}_t$，让它最大化一个\emph{特殊奖励——这个奖励就等于高层的价值函数} $Q^\pi(s, a_t)$，并且只在低层这一段 rollout 的\emph{末尾}给出。

\textbf{一个具体例子（网页智能体）}：状态 = 至今访问过的网页交互；高层动作 = 一句话，如 \texttt{search[queen-sized bed, black]}（搜索“黑色大号床”）；奖励 = 买对了商品给 +1、否则 0；环境动态 = 网页引擎对动作的反应。而这句话里的每个 token（“search”“[”“queen”……）就是\emph{低层 MDP} 里的单个动作。

\subsection{方法预备：RL 的几个定义}

为后面公式做准备，作者复述了标准 RL 定义（用大白话解释）：
\begin{itemize}[leftmargin=1.5em]
\item \textbf{Q 函数} $Q^\pi(s_h, a_h)$：在当前步执行某动作、之后一直按策略 $\pi$ 行动，能拿到的\emph{期望长期回报}。形式上 $Q^\pi(s_h,a_h)=\mathbb{E}_\pi\!\left[\sum_{t=0}^{\infty}\gamma^t r(s_{h+t}, a_{h+t})\right]$，其中 $\gamma$ 是折扣因子（越靠后的奖励越打折）。
\item \textbf{价值函数} $V^\pi(s_h)$：把 $Q$ 在“按策略采样的动作”上取期望，即“处在状态 $s_h$ 的平均好坏”。
\item \textbf{优势函数} $A^\pi(s_h,a_h)=Q^\pi(s_h,a_h)-V^\pi(s_h)$：这个动作\emph{比平均水平好多少}。优势 $>0$ 说明这步比平均强，值得鼓励。
\item 低层 MDP 的价值函数记作 $\widetilde{V}(s_c, a^{1:i-1}_h)$。
\end{itemize}

\subsection{方法核心：在分层 MDP 上跑两套 RL}

分层结构带来一个关键好处——\emph{两层可以各用各的算法}。作者点出二者需求不同：
\begin{quote}
“只有\emph{高层}需要与（不可微的）环境交互；\emph{低层}只针对高层价值函数做优化，因此完全在‘体内（in silico）’训练，不需要任何环境交互。所以两层要求不同：高层算法要\emph{样本效率高}，低层算法要\emph{容易优化}。”
\end{quote}
于是一个特别顺手的选择就是：\textbf{高层用 TD 学习}（off-policy，复用旧数据、样本效率高），\textbf{低层用 on-policy 策略梯度}（好优化，直接继承 RLHF 那一套）。

\subsection{方法落地：面向在线 RL 的高样本效率实例}

这是论文给出的\emph{可直接用}的具体算法，配合图 2 看最清楚（见下一节图解）。它训两个东西：\textbf{话语级 critic}（高层）+ \textbf{token 级 actor}（低层）。

\textbf{高层：话语级 critic}。在高层训\emph{两个} LLM：一个表示话语级 Q 函数 $Q^\pi_\theta(s,a)$，一个表示话语级价值函数 $V^\pi_\psi(s)$。Q 模型用“延迟副本的 V 模型”算出的 Bellman 目标来训；V 模型反过来去逼近“Q 模型在 actor 自回归采样出的句子上的期望值”。因为是 off-policy，所有历史在线交互都存进数据集 $\mathcal{D}=\{s_i,a_i,r_i,s'_i\}_{i=1}^N$ 反复用。两个训练目标（公式 1、2）：

\[ J_Q(\theta)=\mathbb{E}_{s,a,r,s'\sim\mathcal{D}}\left[(Q_\theta(s,a)-r-\gamma V_{\bar{\psi}}(s'))^2\right] \tag{1} \]
\[ J_V(\psi)=\mathbb{E}_{s\sim\mathcal{D}}\left[\mathbb{E}_{a\sim\pi_\varphi(\cdot|s)}\left[(V_\psi(s)-Q_{\bar{\theta}}(s,a))^2\right]\right] \tag{2} \]

\textbf{公式 1 逐符号拆解}：让 $Q_\theta(s,a)$ 去拟合“即时奖励 $r$ + 折扣 $\gamma$ × 下一状态的价值 $V_{\bar\psi}(s')$”。这就是 Bellman 一致性：一个动作的价值 = 它当下拿的奖励 + 它把你带到的下个局面的价值。$\bar\psi$ 上面那条横杠表示\emph{延迟目标网络}（target network，缓慢更新的副本），用来稳住训练、防止自己追自己发散。注意——\emph{这里的 $s,a,s'$ 都是话语级的}，所以 Bellman 回溯是\emph{一句话跨到下一句话}，而不是一个 token 跨到下一个 token，时间跨度被大大缩短。

\textbf{公式 2 逐符号拆解}：让 $V_\psi(s)$ 去逼近“在状态 $s$ 下，按当前 actor $\pi_\varphi$ 采样动作 $a$，再用目标 Q 模型 $Q_{\bar\theta}$ 打分”的期望。\emph{巧妙之处}：本来话语级 TD 学习需要“在所有可能句子里取最大值”（不可行），这里改成“\emph{从 actor 采样几句话再求期望}”，用 actor 产生的真实句子来近似，绕开了对天文数字动作空间求 max。两个目标网络 $Q_{\bar\theta}, V_{\bar\psi}$ 用 Polyak 平均（缓慢滑动平均）更新。

\textbf{低层：token 级 actor}。用 on-policy 策略梯度训 token 级策略 $\pi_\varphi(\cdot|s_c, a^{1:h}_t)$，去找“能最大化 Q 模型预测”的 token 序列。为降低方差，用从 Q 模型导出的\emph{优势值}当作\emph{末尾奖励}。这个子过程其实就是\emph{把 RLHF 的单轮 RL 推广了一下，只不过末尾奖励从“奖励模型”换成了“多轮优势的估计”}。用 REINFORCE 更新（公式 3）：

\[ J_\varphi(\pi)=\mathbb{E}_{s_c\sim\mathcal{D},\,a^{1:L}_t\sim\pi(\cdot|s_c)}\left[\sum_{i=1}^{L} A(s_c, a^{1:L}_t)\,\log\pi_\varphi(a^i_t|s_c, a^{1:i-1}_t)\right] \tag{3} \]

\textbf{公式 3 逐符号拆解}：对生成的每个 token $a^i_t$，用“整句话的优势 $A(s_c, a^{1:L}_t)$”去加权它的对数概率 $\log\pi_\varphi(\cdot)$ 的梯度。优势为正（这句话比平均好）就提高这串 token 的概率，为负就压低。注意优势在\emph{句末}才算出，再分摊到本句每个 token——这正是图 2 里“把优势 $A^\pi(s_c,a_t)$ 当作 token 序列末尾奖励”的含义。

\subsection{方法变体：还能怎么改}

ArCHer 是个\emph{框架}，作者给了两个有用变体：

\textbf{(1) 改进 token 级策略梯度（降方差）}。当每句话很长、token 级时间跨度长时，公式 3 的 REINFORCE 方差很大、难改进策略。办法是在 token 级 MDP 里引入一个\emph{基线价值函数} $\widetilde{V}_\eta$（用蒙特卡洛回报做监督回归来训，公式 4），再把它从优势里减掉（公式 5）：

\[ J_\eta(\widetilde{V})=\mathbb{E}_{s_c\sim\mathcal{D},\,a^{1:L}_t\sim\pi(\cdot|s_c)}\left[\sum_{i=1}^{L}\left(A(s_c,a^{1:L}_t)-\widetilde{V}_\eta(s_c,a^{1:i-1}_t)\right)^2\right] \tag{4} \]
\[ J_\varphi(\pi)=\mathbb{E}_{s_c\sim\mathcal{D},\,a^{1:L}_t\sim\pi(\cdot|s_c)}\left[\sum_{i=1}^{L}\left(A(s_c,a^{1:L}_t)-\widetilde{V}_\eta(s_c,a^{1:i-1}_t)\right)\log\pi_\varphi(a^i_t|s_c,a^{1:i-1}_t)\right] \tag{5} \]

直觉：基线 $\widetilde{V}_\eta$ 像一条“及格线”，只奖励/惩罚\emph{超出或低于及格线的部分}，让信号更干净。代价是要多训一个网络（$V_\eta$ 用单独的 GPT-2 参数化），句子短时未必划算。

\textbf{(2) 离线 ArCHer（完全不在线交互）}。直接拿一批历史数据学。离线设定的麻烦是：改进策略往往会选到\emph{分布外动作}（out-of-distribution, OOD），而这些动作的 $Q$ 值在离线数据里估不准，硬优化会\emph{严重高估、发散、性能崩}。解法是换目标函数：用 \textbf{IQL}（implicit Q-learning，隐式 Q 学习）给 critic 算 Bellman 目标（只在“数据支持的动作”上做），用 \textbf{AWR}（advantage-weighted regression，优势加权回归）约束 actor 别离数据太远。

IQL 损失（公式 6）：用一个\emph{期望分位数（expectile）回归}来\emph{隐式}地近似“对动作取最大值”，而不用真的去 max：
\[ J^{\mathrm{IQL}}_\psi(V)=\mathbb{E}_{s\sim\mathcal{D}}\left[\mathbb{E}_{a\sim\pi_\varphi(\cdot|s)}\left[L^\tau_2(V_\psi(s)-Q_{\bar\theta}(s,a))\right]\right] \tag{6} \]
其中 $L^\tau_2(u)=|\tau-\mathbf{1}\{u<0\}|\,u^2$，$\tau\in[0.5,1)$ 是期望分位数参数（$\tau$ 越大越偏向“乐观的高值”，即越接近取 max）。

AWR 损失（公式 7）：在“追高回报”和“模仿生成数据的策略”之间折中：
\[ J_\varphi(\pi)=-\mathbb{E}_{(s_c,a^{1:L}_t)\sim\mathcal{D}}\left[\exp\!\left(\beta\cdot A(s_c,a^{1:L}_t)\right)\sum_{i=1}^{L}\log\pi_\varphi(a^i_t|s_c,a^{1:i-1}_t)\right] \tag{7} \]
$\beta$ 是用户设的正标量：$\beta$ 小 → 偏向模仿原数据（稳）；$\beta$ 大 → 更激进地追奖励（可能不稳）。

\subsection{框架总结与实现细节}

\textbf{伪代码（原文 Algorithm 1）的脉络}：每轮迭代里——(a) 在线模式下用当前 actor 收集数据进缓冲区；(b) 用 Bellman 自举更新话语级 Q、V（公式 1 与 公式 2 或 6），并用滑动平均更新目标网络；(c)（可选）用蒙特卡洛回归更新 token 级基线（公式 4）；(d) 用话语级 critic 给的优势更新 token 级 actor（公式 3、5 或 7）。同一套伪代码既能跑\emph{在线}也能跑\emph{离线}，且高/低层目标可灵活替换——这正是“框架”的意义。

\textbf{关键实现选择}：策略（actor）用 \textbf{GPT-2} 架构；critic 用 \textbf{RoBERTa-base}（取 “[CLS]” token 的嵌入接一层线性层出预测）。为对抗 Q 值高估，用\emph{双 Q 学习}（double Q-learning）trick：独立训两套 $\{Q_1,V_1\}$、$\{Q_2,V_2\}$，优势用 $Q_1,Q_2$ 取最小、$V_1,V_2$ 取最小来算（取 min 更保守，防高估）。为省算力/显存，四个网络\emph{共享同一个语言模型编码器骨干}，只用各自独立的 MLP 头；actor 的参数与 critic 独立。

\subsection{理论分析：为什么“分层”在理论上也更优}

作者证明了分层设计相比 token 级 off-policy 方法，\emph{对算法收敛的条件要求更弱}，且\emph{有限样本下的统计误差更小}。两点直觉（对应 Lemma 1、2）：

\begin{itemize}[leftmargin=1.5em]
\item \textbf{Bellman 完备性（函数类容量）}：要满足 Bellman 完备这个条件，\emph{话语级} critic 的函数类要求\emph{更弱}。直觉是 token 级 critic 必须能在“下一个 token”这个细粒度上拟合任意函数（要求容量很高），而话语级 critic 只需在“话语”这个更粗的时间尺度上拟合任意函数。
\item \textbf{数据覆盖（密度比条件）}：话语级和 token 级\emph{要求一样}。直觉是同一批离线数据在两层诱导出的轨迹分布是一样的——一条轨迹在话语级被覆盖，在 token 级也被覆盖，反之亦然。
\end{itemize}

\textbf{主定理（Theorem 1，非正式）的结论}：在最坏情况下，\emph{用 token 级 critic 估计优势的误差，比用话语级 critic 大约大} $\sqrt{\gamma^L}$ 倍（$L$ 是每句话最多的 token 数），原因是\emph{误差累积}。实践里 $\gamma$ 常 $>0.95$、$L$ 可达 64，导致 $\gamma^{L^{1/2}}$ 这个因子远大于 1。换句话说：分层设计不仅收敛条件更弱，统计误差保证也更好，最终给出的策略梯度估计在最坏情况下更准。

\section{方法图解：ArCHer 的架构与数据流}

图~\ref{fig:arch} 是 ArCHer 实用版的整体示意（原文 Figure 2），上下分两层，是理解全文的关键。

\begin{figure}[H]\centering
  \includegraphics[width=0.96\linewidth]{cand_p04_k00.pdf}
  \caption{ArCHer（带分层结构的演员-评论家框架）实用实例示意。算法同时在\emph{话语级}与 \emph{token 级}运作。在话语级，通过带 TD 误差的 Bellman 自举学习一个 Q 函数；在 token 级，用策略梯度法最大化“由话语级 Q 函数诱导的优势函数”，并把这个优势估计作为本句 token 序列\emph{末尾}的奖励。（图内英文为原文保留。）}
  \label{fig:arch}
\end{figure}

\textbf{怎么读这张图}：

\begin{itemize}[leftmargin=1.5em]
\item \textbf{上半部分（Utterance Level，话语级）}：把\emph{一整句话}当一个动作。左边的 Utterance-Level Critic 吃进“状态 $s$ + 动作 $a$（一句话）”，输出 $Q^\pi(s,a)$；右边同一个 critic 吃进“下一状态 $s'$”，输出 $V^\pi(s')$。中间橙色箭头 \textbf{Bellman Backup}（Bellman 回溯）执行公式 1：$Q^\pi(s,a)\leftarrow r+\gamma V^\pi(s')$——\emph{注意箭头是一句话跨到下一句话}，所以时间跨度短。这一层用 off-policy 数据训，样本效率高。图中那个“买黑色大号床”的例子展示了状态就是“与购物网站来回的交互历史”，动作就是 \texttt{search[black queen-sized bed]} 这样的一句话。
\item \textbf{下半部分（Token Level，token 级）}：把\emph{这句话里的每个字}当一个动作。Token-Level Actor 在“本轮上下文 $s_c$”基础上，一个字一个字地自回归生成 $a^1_t, \dots, a^L_t$（如 “Search”“[”“black”……“]”“EOS”）。
\item \textbf{两层的连接（橙黄色 Policy Gradient 大箭头）}：话语级 critic 算出的\emph{优势} $A^\pi(s_c, a_t)$，被当作\emph{这句话末尾的奖励}，顺着大箭头“喂”给下层的 token 级 actor，用策略梯度（公式 3/5）更新它。\emph{这就是“高层评委打分、低层据此学怎么说话”的数据流}。
\end{itemize}

一句话概括这张图：\emph{高层（话语级）当评委、用能复用旧数据的 TD 学习快速学会“评估一句话的好坏”；低层（token 级）当演员、用 PPO 那套 on-policy 梯度、拿评委给的优势分当奖励，学会“具体每个字怎么生成”。}

\section{实验解读}

\subsection{用了哪些任务、比了哪些基线}

\textbf{评测任务（5 个）}。作者特意挑了\emph{需要多步规划、不能一轮解决、7B 内模型能做、且评测快而可复现}的任务：
\begin{itemize}[leftmargin=1.5em]
\item \textbf{Detective Game}（侦探游戏）：交互式文字冒险，要靠环境反馈一步步生成自然语言动作（如“take paper”），最终目标是找出凶手；到达里程碑给奖励。
\item \textbf{Twenty Questions}（二十问）：智能体当“猜词者”，在 20 个是/否问题内从 157 个词里猜出隐藏词。神谕（oracle）用 flan-t5-small 模拟。猜对得 0、终止；否则每步 $-1$。还有一个只含 10 个词的简化版 \textbf{Twenty Questions Subset}（考“走捷径”的专门策略）。
\item \textbf{Guess My City}（猜城市）：从 100 个城市里 20 问内猜出隐藏城市。与二十问的关键区别：可以问\emph{任意问题}、答案是\emph{自由文本}而非只有是/否（更长更难）。
\item \textbf{WebShop}：工具使用任务，按指令去购物服务器上买东西，按“买到的商品与要求的相似度”给 0$\sim$1 的稠密奖励。
\end{itemize}

\textbf{对比基线}：
\begin{itemize}[leftmargin=1.5em]
\item \textbf{token 级 PPO}：当前 LLM 上的 SOTA on-policy 方法，但旧数据用完即弃。
\item \textbf{Filtered BC}（过滤式行为克隆）：非 RL 基线，只模仿最近 rollout 中\emph{回报最高的前 10\%}。作者强调：这个简单基线常常出奇地强，所以“能否超过 Filtered BC”是判断一个方法里“RL 成分”是否真正起作用的\emph{试金石}。
\item \textbf{CHAI}：SOTA 话语级 RL 方法。关键区别——CHAI 用一个\emph{冻结的} actor（行为克隆得到）采样 $k$ 个候选句子，再让 critic 挑 Q 值最高的；而 ArCHer\emph{会一起优化 actor}。实验取 $k=5$（即便如此，CHAI 运行时间已是 ArCHer 的约 4 倍）。
\item 在 WebShop 上还比了 \textbf{GPT-3.5 + 专家提示} 和 \textbf{GPT-3.5 + ReAct}。
\end{itemize}
所有方法的 token 级策略都用“在次优数据上做监督指令微调”的同一个 checkpoint 初始化，保证公平、且利于 RL 初期探索。

\subsection{主要结果：100 倍样本效率 + 更强最终性能}

\begin{table}[H]\centering\caption{核心结论速览（依据原文图 3、图 4、5.3 节文字整理）。}
\begin{tabular}{@{}p{4.2cm}p{10.5cm}@{}}
\toprule
\textbf{对比对象} & \textbf{结论} \\ \midrule
vs.\ PPO（on-policy） & 在 4 个较难任务上 ArCHer 收敛到\emph{明显更好}的性能；样本效率约 \textbf{100 倍}：二十问上 PPO 要 $>$10 万样本才让平均回报勉强高于 $-17$，ArCHer 不到 \textbf{1000} 样本就达到同一水平。 \\
vs.\ Filtered BC & Filtered BC 收敛很快但\emph{停在次优}、给更多数据也不再提升；ArCHer 随样本增多\emph{持续稳定改进}。 \\
vs.\ CHAI（话语级 RL） & CHAI 初期比 ArCHer 快，但\emph{最终收敛得更差}（因为它只在冻结策略上重排候选，改进空间窄）；ArCHer 需要一段初期学习把 critic 误差降下来，之后稳步超越。 \\
WebShop & 用 ArCHer 微调的 \emph{GPT-2}，是\emph{唯一}超过“GPT-3.5 + 专家提示 / ReAct”的方法。 \\
\bottomrule
\end{tabular}
\end{table}

整体上：ArCHer 在 4 个更难的任务（Twenty Questions Subset、Twenty Questions、Guess My City、WebShop）上\emph{收敛到比所有前人方法都好}的性能；在最简单的 Detective Game 上则与最佳前人方法\emph{打平}。

\subsection{消融实验：拆开看每个设计是否有用}

\begin{itemize}[leftmargin=1.5em]
\item \textbf{off-policy 数据的重要性（图 5b，Guess My City）}：把回放缓冲区调小（只留最近 48 条 rollout）→ 学习\emph{不稳定}（在有限数据上过拟合）；缓冲区\emph{大一些更稳}，但大到一定程度后再增大就没明显收益。结论：用好 off-policy 数据能提升稳定性与性能。
\item \textbf{critic 用什么骨干（图 5c，Twenty Questions）}：把 critic 从\emph{仅编码器}的 RoBERTa 换成\emph{仅解码器自回归}的 GPT-2（取最后 “[EOS]” token 嵌入）。结果：RoBERTa 版开头略快，但过了初期两者\emph{表现一致}。结论：ArCHer 用解码器模型当 critic 也行，不掉性能。
\item \textbf{模型规模 100M→7B（图 5d，Twenty Questions Subset）}：把 token 级 actor 从 1 亿参数的 GPT-2 换成 70 亿的 \textbf{Mistral-7B}（7B 开源 checkpoint 零样本已有不俗表现，无需再做监督微调）。结果：Mistral-7B 版\emph{学得快得多}。结论：ArCHer 能随模型变大而\emph{更好地 scale}——这与“actor 更新本质上和 RLHF 单轮微调相似、应同样受益于规模”的预期一致。
\item \textbf{token 级基线（图 5a，Guess My City）}：这个任务句子长而多样。加上 token 级基线（公式 5）\emph{大幅超过}不加基线的标准 ASCHer，验证了“基线能在长/多样句子时有效降方差”的假设；代价是要多训一个 $\widetilde{V}_\eta$，句子短时未必值。
\end{itemize}

\subsection{离线设定：IQL + AWR 最稳（表 1）}

\begin{table}[H]\centering\caption{ArCHer 在\emph{离线} RL 设定下的不同变体（依据原文 Table 1，Twenty Questions 上的回报 Return，跑 1280 条轨迹 / 5 个随机种子）。}
\begin{tabular}{@{}lc@{}}
\toprule
\textbf{方法（离线）} & \textbf{Twenty Questions（Return）} \\ \midrule
ArCHer (IQL + AWR)            & $-14.1$ \\
ArCHer (IQL + REINFORCE)      & $-20$ \\
ArCHer (IQL + REINFORCE + BC) & $-15.3$ \\
ArCHer (SARSA + AWR)          & $-14.5$ \\
Filtered BC                   & $-15.4$ \\
BC                            & $-16.8$ \\
\bottomrule
\end{tabular}
\end{table}

读这张表（回报越接近 0 越好）：
\begin{itemize}[leftmargin=1.5em]
\item 直接把在线的 REINFORCE（公式 3）搬到离线 → \emph{性能迅速崩塌}（$-20$），因为没有任何约束防分布外动作，这是离线 RL 的经典坑。
\item 给它加上模仿损失（+BC）能稳住、提升到 $-15.3$，但仍\emph{不如 AWR}。
\item 把 IQL 的“数据内分位数回溯”换成 \textbf{SARSA} 回溯（直接用数据里下一句话算 Bellman 目标、\emph{不做任何取最大值}）→ $-14.5$，不如 IQL（$-14.1$）。这\emph{凸显了离线设定下“对动作取最大值”的重要性}。
\item \textbf{结论}：用 IQL + AWR（best，$-14.1$）这种专门处理分布外动作的组合最好；且 IQL/SARSA 配 AWR 都\emph{优于}朴素 BC 和 Filtered BC，说明“动态规划训 critic”确实有用。
\end{itemize}

\section{核心 takeaways（给 AI 应用开发者）}

\begin{enumerate}[leftmargin=1.6em]
\item \textbf{“分层”是处理多轮长程任务的关键招式}。当你的任务是“模型要来回好几轮才能完成”（客服、网页操作、工具链、对话搜集需求），不要硬套单轮 RLHF。ArCHer 的思路——\emph{高层管“说哪句话”、低层管“这句话怎么生成”}——把“超长时间跨度”和“超大动作空间”这对矛盾一次性化解。

\item \textbf{样本效率 = 真金白银}。多轮任务里每次与环境交互（真去点网页、真调 API、真等用户回复）都慢且贵。能复用历史数据的 off-policy 设计带来约 \emph{100 倍}样本效率（相对 PPO），意味着同样的预算下少跑两个数量级的交互。做 agent 训练时，\emph{优先考虑能复用 replay buffer 的方案}。

\item \textbf{“能否打败 Filtered BC”是检验 RL 是否真起作用的试金石}。过滤式行为克隆（只模仿历史最好的前 10\%）简单到几乎是免费基线，却常常很强。你自研的 RL 方案如果\emph{打不过它}，多半是“RL 成分”没真正生效——这是个极实用的 sanity check。

\item \textbf{在线 / 离线要用不同的“配方”}。在线可用简单的 REINFORCE；但一旦改成\emph{从静态历史数据离线学}，必须换成能压制\emph{分布外动作}的目标（IQL 给 critic、AWR 给 actor），否则会高估 Q 值、训练发散。把在线算法原样搬到离线是常见翻车点。

\item \textbf{这套框架对模型/算法是“即插即用”的}。actor、critic 的骨干可换（GPT-2、RoBERTa、Mistral-7B 都行），高/低层算法可换（TD / IQL / SARSA、REINFORCE / AWR）。而且 actor 的更新本质上就是 RLHF 那一套，\emph{能直接复用现有单轮 RL 的工程实现}，并随模型变大而 scale。

\item \textbf{已知边界（用前要知道）}。论文主实验多用较小的 GPT-2；虽有 Mistral-7B 的良好 scaling 迹象，但\emph{大模型 + 更多基准}的系统验证仍是未来工作。此外方法仍需\emph{数千次}交互、且只在“可计算的奖励”任务上验证——若你的场景只有\emph{约 100 次}真人交互的预算，作者认为基于模型的 RL（model-based RL）可能更有前途。
\end{enumerate}

\section*{参考文献（仅列与本篇最相关的几条，英文原样保留）}
\begin{enumerate}[leftmargin=2em]
\item Yifei Zhou, Andrea Zanette, Jiayi Pan, Sergey Levine, Aviral Kumar. ArCHer: Training Language Model Agents via Hierarchical Multi-Turn RL. ICML, 2024. \url{https://github.com/YifeiZhou02/ArCHer}
\item Charlie Snell et al. Offline RL for Natural Language Generation with Implicit Language Q Learning (ILQL). ICLR, 2023.
\item Siddharth Verma et al. CHAI: A Chatbot AI for Task-Oriented Dialogue with Offline RL. NAACL, 2022.
\item Ilya Kostrikov, Ashvin Nair, Sergey Levine. Offline Reinforcement Learning with Implicit Q-Learning (IQL). ICLR, 2022.
\item Xue Bin Peng et al. Advantage-Weighted Regression (AWR): Simple and Scalable Off-Policy Reinforcement Learning, 2019.
\item John Schulman et al. Proximal Policy Optimization Algorithms (PPO). arXiv:1707.06347, 2017.
\item Marwa Abdulhai et al. LMRL Gym: Benchmarks for Multi-Turn Reinforcement Learning with Language Models, 2023.
\end{enumerate}

\end{document}
